% Contenido del póster: Kim y Stähler (2026), Journal of Economics and Business. DOI 10.1016/j.jeconbus.2026.106307
% Elaborado con apoyo de IA a partir del resumen oficial del artículo (OpenAlex), 2026-11-12.
\def\DOI{10.1016/j.jeconbus.2026.106307}
\def\TituloEs{El impacto de los préstamos entre pares en el crédito a pequeñas empresas}
\def\TituloOriginal{The impact of peer-to-peer lending on small business loans}
\def\Autores{Kim y Stähler}
\def\Anio{2026}
\def\Revista{Journal of Economics and Business}

\def\Problema{Las plataformas de préstamos entre pares compiten con los bancos por los mismos prestatarios. No está claro si su llegada reemplaza el crédito bancario a las pequeñas empresas o si amplía el crédito total disponible.}
\def\Objetivo{Medir el efecto de los préstamos entre pares en los créditos a pequeñas empresas que otorgan las instituciones de depósito de Estados Unidos sujetas a la Community Reinvestment Act.}
\def\Metodo{Los autores plantean un modelo en el que el prestatario elige entre un banco tradicional y una plataforma de crowdlending. Luego aprovechan la entrada escalonada de LendingClub en los estados de Estados Unidos entre 2009 y 2017 para estimar cómo cambió el crédito bancario a pequeñas empresas en cada zona.}
\def\Resultados{\begin{itemize}
  \item En el modelo, la entrada de la plataforma genera dos efectos: un cambio de prestatarios del banco a la plataforma y una expansión del crédito.
  \item Con la entrada de LendingClub bajaron los créditos bancarios a pequeñas empresas.
  \item La caída fue mayor en las zonas de ingreso bajo o moderado y en las zonas de ingreso medio con alta pobreza.
  \item Una estimación conservadora sugiere que el volumen total de crédito a pequeñas empresas pudo disminuir.
\end{itemize}}
\def\Conclusion{En este caso, el crédito de plataforma actuó más como sustituto que como complemento del crédito bancario, y la sustitución se concentró en las zonas más vulnerables.}
\def\Relevancia{En el Perú, el financiamiento participativo financiero se volvió casi exclusivamente empresarial: en 2026 los proyectos empresariales explican el \cifra{partEmp2026}~\% del monto. Para saber si sustituye al crédito bancario a la pequeña y microempresa haría falta el volumen de ese crédito por zona, que no está en las series mensuales usadas. En el artículo del autor, los cambios del monto financiado no se asocian con los cambios de las tasas tradicionales.}
\def\PalabrasClave{préstamos entre pares, crédito a pequeñas empresas, sustitución, LendingClub, Community Reinvestment Act}
% --- Ficha del estudio (herramientas/posteres/ficha_articulos.py) ---
\def\Ambito{Estados Unidos (LendingClub, 2009-2017)}
\def\Enfoque{Modelo teórico y entrada escalonada de la plataforma}
\def\FiguraPeru{3}
% --- fin de la ficha ---
